\chapter{The Tower}
\label{chap:the-tower}

% EPISODE COVERAGE: Episodes 5, 6, and 7 -- SOURCE/EXECUTED/Abstraction,
% VALUE/MAGNITUDE/SCALED/LOAD/FINITE_ELEPHANT/Spline, Episode 7 totality proof
% KEY CONCEPT: The totality proof. The tower exists not just as definitions
% but as a concrete verified instance.

\section{Encoding Instructions}
\label{se:encoding-instructions}

% Episode 5: Equivalation, SOURCE, Encoding, EXECUTED, Abstraction.
%
%   inductive Encoding where
%     | boot: Encoding
%     | zero: Encoding
%     | one: Encoding
%
% Three opcodes: boot (initialize), zero, one. Every program is a
% sequence of these. This is the minimal instruction set.
% EXECUTED wraps this with execute? and output?.
%
% Sidebar: The instruction set as a measuring instrument. Every CPU
% has an instruction set -- the finite vocabulary of operations it can perform.
% The Encoding type is the minimal instruction set: boot, zero, one.
% Everything else is derived.

\section{Abstraction and the Berry Paradox}
\label{se:abstraction-berry}

% Abstraction:
%   inductive Abstraction where
%     | satire: Abstraction
%     | compile: Abstraction
%     | execute: Abstraction
%
% Abstraction.le and Abstraction.lt: careful ordering to avoid the Berry paradox.
% The Berry paradox arises from self-referential definition. Abstraction avoids
% this by requiring that lt terminate -- compile < execute, satire incomparable.
%
% VALUE wraps this with monad (the abstraction of the closure) and lt?.
%
% Sidebar: The Berry paradox and self-reference. Russell, Godel, Chaitin,
% and the halting problem are all variations: what happens when a formal
% system refers to itself? Abstraction navigates this by making the
% reference structure explicit and bounded.

\begin{gphenom}{Galileo-Abel}
\label{ph:galileo-abel}
\GPhObs Galileo showed that all objects fall at the same rate regardless of mass.
Abel showed that the roots of a polynomial are invariant under permutation of
the coefficients. In both cases, the law is preserved under a transformation
that the experiment does not see.
\GPhCon If two instruments disagree on a measurement but agree on the law, one
of them is applying a transformation the law does not see.
The law is the invariant; the instrument choice is the gauge.
\GPhSeq \lean{VALUE}'s ordering on \lean{Abstraction} is invariant under
transformation: \lean{compile < execute} regardless of which instrument
is doing the compiling. The law holds across instruments.
\end{gphenom}

\section{Arithmetic Assembled}
\label{se:arithmetic-assembled}

% Episode 6: Sum, MAGNITUDE, Product, SCALED, Basis, LOAD, Polynomial,
% FINITE_ELEPHANT, Spline.
%
% Sum: zero and add. Addition.
% MAGNITUDE: adds whelmed? -- can one Sum dominate another?
% Product: origin (vacuum), one (identity), mul (scaling).
% SCALED: adds orthogonal? -- can two Products be orthogonal?
% Basis: null_space, origin, basis -- the three-part vector space structure.
% LOAD: adds decoded? -- can we invert the encoding?
% FINITE_ELEPHANT: the DAC layer. Takes discrete values, produces smooth oscillations.
% BASICProcess: GOSUB/GOTO. The compiler's minimal flow control.
%
% Sidebar: Why BASIC? The minimal model of flow control requires unconditional jumps.

\begin{gphenom}{Dirichlet-Bancroft}
\label{ph:dirichlet-bancroft}
\GPhObs Dirichlet's function is 1 on rationals and 0 on irrationals.
It is defined everywhere but Riemann-integrable nowhere.
A model can be complete -- defined for every input -- while being
useless for measurement.
\GPhCon Model completeness is not ledger sufficiency. A model that cannot
invert its encoding cannot recover the original record, even if it
is defined for every possible input.
\GPhSeq \lean{LOAD}'s \lean{decoded?} predicate marks this boundary.
A model accepted by the compiler is not necessarily a model that
can answer questions. Completeness is cheap; sufficiency is earned.
\end{gphenom}

\section{The Spline}
\label{se:the-spline}

% Spline:
%   inductive Spline where
%     | observation: Spline
%     | knot: Spline
%     | interpolant: Spline
%
% The continuum made explicit. The spline is not assumed -- it is constructed
% from observations and knot constraints. The smooth curve is the unique
% admissible interpolant: minimum curvature, threading all observed points.
%
% Polynomial and GalerkinProcess: the algebraic engine. The Galerkin method
% projects the continuous problem onto a finite polynomial basis.
% The compiler is doing finite element analysis.

\section{The Totality Proof}
\label{se:totality-proof}

% Episode 7: 17 chained instance declarations.
% Each instance is the compiler confirming that the previous level was real.
% To declare an instance is to produce a concrete inhabitant of the typeclass.
%
% The cascade is the totality proof: DISTINGUISHABLE through ACOLYTE, verified.
% The tower is not a collection of definitions that might be consistent.
% It is a construction that has been built and checked.
%
% The heartbeat counter: maxHeartbeats must be raised here. The compiler sweats.

\begin{gphenom}{Markov-Conway}
\label{ph:markov-Conway}
\GPhObs Conway's Game of Life produces arbitrarily complex structures --
gliders, oscillators, universal computers -- from three local rules
applied to a binary grid. No global architect is required.
\GPhCon Global structure emerges from local admissibility conditions.
The tower does not need a designer; it needs each floor to be
locally consistent with the floor below.
\GPhSeq Each typeclass instance in \episode{7} is a local rule.
The 17-instance cascade is the Game of Life running on the tower:
global structure -- the verified existence of the entire hierarchy --
emerging from 17 local consistency checks.
\end{gphenom}

% End of chapter: the tower stands. Seventeen floors, each verified.
% The next floor asks something harder: what is true?

